%=====================================================================
\chapter{Sequences, Time Series and Streaming Analytics}\label{ch:timeseries}
%=====================================================================
\locator{4}{Part IV --- Deep Learning and Modern AI}

\begin{outcomes}
By the end of this chapter you will be able to:
\begin{tight}
  \item \textbf{Decompose} a time series into trend, seasonality and residual, and read the decomposition.
  \item \textbf{Build} lag, rolling-window and calendar features for a forecasting model.
  \item \textbf{Validate} a temporal model correctly, and explain why random splitting is invalid.
  \item \textbf{Distinguish} batch from stream processing and describe windowing.
  \item \textbf{Detect} concept drift and specify a response.
\end{tight}
\end{outcomes}

\begin{prereq}
\chref{descriptive}, \chref{evaluation} (validation strategy) and
\chref{features} (window aggregations). This chapter is the direct bridge into
\chref{iot}.
\end{prereq}

\begin{keyterms}
time series $\bullet$ trend $\bullet$ seasonality $\bullet$ residual $\bullet$
stationarity $\bullet$ lag feature $\bullet$ rolling window $\bullet$
autocorrelation $\bullet$ ARIMA $\bullet$ horizon $\bullet$ walk-forward validation
$\bullet$ batch vs stream $\bullet$ tumbling and sliding windows $\bullet$
watermark $\bullet$ concept drift
\end{keyterms}

\section{What Makes Time Different}

\begin{alertbox}[The Rule That Governs This Whole Chapter]
Observations are \textbf{not independent}. Today's value depends on yesterday's.
Every technique in Part III assumed independence, so every one of them needs
adjusting here --- most importantly the validation strategy. Shuffling a time
series destroys the only structure that matters.
\end{alertbox}

\section{Decomposition}

\begin{definitionbox}[Time Series Decomposition]
Any series can be viewed as
\[
y_t = \text{Trend}_t + \text{Seasonality}_t + \text{Residual}_t
\]
(or multiplicatively). \term{Trend} is the long-run direction, \term{seasonality}
is the repeating cycle, and the \term{residual} is what is left --- which is where
anomalies live.
\end{definitionbox}

\dsfig{%
\begin{tikzpicture}
\begin{groupplot}[
  group style={group size=1 by 4, vertical sep=0.42cm, x descriptions at=edge bottom},
  width=12.4cm, height=2.5cm,
  axis lines=left, tick label style={font=\tiny}, label style={font=\tiny},
  title style={font=\tiny\bfseries, at={(0.02,0.92)}, anchor=north west},
  xmin=0, xmax=48, ytick=\empty,
]
\nextgroupplot[title={\color{primary}OBSERVED --- weekly VLE logins}]
\addplot[primary, line width=0.9pt, domain=0:48, samples=145]
  {40 + 0.9*x + 12*sin(deg(2*pi*x/12)) + 4*sin(deg(x*3.1))};

\nextgroupplot[title={\color{secondary}TREND --- steady growth}]
\addplot[secondary, line width=1.1pt, domain=0:48, samples=100] {40 + 0.9*x};

\nextgroupplot[title={\color{greenhl}SEASONALITY --- a 12-week teaching cycle}]
\addplot[greenhl, line width=0.9pt, domain=0:48, samples=145] {12*sin(deg(2*pi*x/12))};

\nextgroupplot[title={\color{accent}RESIDUAL --- anomalies live here},
               xlabel={week}]
\addplot[accent, line width=0.7pt, domain=0:48, samples=145] {4*sin(deg(x*3.1))};
\addplot[accent, only marks, mark=*, mark size=2pt] coordinates {(31,-13)};
\node[font=\tiny, accent, anchor=west] at (axis cs:32,-12) {outage};
\end{groupplot}
\end{tikzpicture}}%
{Decomposition of a weekly series. Separating trend and seasonality is what lets you say
whether a drop is alarming: a fall in week 31 means nothing if week 31 always falls, and
everything if it does not.}{decomposition}

\begin{conceptbox}[Why Decomposition Is the First Step]
``Logins fell 30\% this week'' is not information until you know whether they fall
every year at this point. Decomposition converts a raw number into a statement
about the residual --- and only the residual is worth alerting on. This is the
single most useful idea in operational monitoring.
\end{conceptbox}

\section{Features for Forecasting}

\rowcolors{2}{white}{lightbg}
\begin{longtable}{@{}p{3.1cm}p{4.6cm}p{6.5cm}@{}}
\toprule
\rowcolor{primary!15}
\textbf{Feature type} & \textbf{Construction} & \textbf{Captures} \\
\midrule
\endhead
Lags & $y_{t-1}$, $y_{t-7}$, $y_{t-365}$ & Recent level, last week, last year --- pick lags matching known cycles \\
Rolling statistics & mean, sd, min, max over a window & Local level and volatility \\
Rolling differences & $y_t - y_{t-1}$; $y_t / y_{t-7}$ & Change and growth rate rather than level \\
Calendar & hour, day of week, month, is\_holiday & Human rhythms --- essential in security and IoT \\
Expanding statistics & mean since the start & Long-run baseline for that entity \\
Time since event & days since last failure / login & Recency, which is often the strongest single feature \\
\bottomrule
\end{longtable}

\begin{alertbox}[Every Lag Must Look Backwards Only]
A rolling mean centred on $t$ uses $y_{t+1}$, which will not exist at prediction
time. Use \emph{trailing} windows exclusively. Most time-series leakage is a
centred window or a \texttt{shift()} in the wrong direction --- and it produces
spectacular, entirely fake, validation scores.
\end{alertbox}

\section{Validating a Temporal Model}

\dsfig{%
\begin{tikzpicture}[font=\scriptsize]
% ---- WRONG
\node[font=\scriptsize\bfseries, text=accent, anchor=west] at (-0.8,1.5)
     {\faIcon{times-circle}~WRONG --- random k-fold: trains on the future to predict the past};
\foreach \i in {0,...,11}{
  \pgfmathsetmacro{\xx}{\i*0.95}
  \pgfmathtruncatemacro{\m}{mod(\i*5,3)}
  \ifnum\m=1
    \node[rounded corners=1.5pt, fill=goldhl!35, draw=goldhl, minimum width=0.85cm,
          minimum height=0.42cm] at (\xx,0.85) {};
  \else
    \node[rounded corners=1.5pt, fill=primary!14, draw=primary!45, minimum width=0.85cm,
          minimum height=0.42cm] at (\xx,0.85) {};
  \fi}
\node[font=\tiny, text=gray!55!black, anchor=west] at (11.4,0.85) {time $\rightarrow$};

% ---- RIGHT
\node[font=\scriptsize\bfseries, text=greenhl!50!black, anchor=west] at (-0.8,-0.15)
     {\faIcon{check-circle}~RIGHT --- walk-forward: always train on the past, test on the future};
\foreach \r/\trainend in {1/3, 2/5, 3/7, 4/9}{
  \pgfmathsetmacro{\yy}{-0.75-(\r-1)*0.6}
  \foreach \i in {0,...,11}{
    \pgfmathsetmacro{\xx}{\i*0.95}
    \ifnum\i<\trainend
      \node[rounded corners=1.5pt, fill=primary!14, draw=primary!45, minimum width=0.85cm,
            minimum height=0.42cm] at (\xx,\yy) {};
    \else
      \ifnum\i<\numexpr\trainend+2\relax
        \node[rounded corners=1.5pt, fill=goldhl!35, draw=goldhl, minimum width=0.85cm,
              minimum height=0.42cm] at (\xx,\yy) {};
      \else
        \node[rounded corners=1.5pt, fill=gray!8, draw=gray!25, minimum width=0.85cm,
              minimum height=0.42cm] at (\xx,\yy) {};
      \fi
    \fi}
  \node[font=\tiny, text=gray!55!black, anchor=west] at (11.4,\yy) {split \r};}

\node[anchor=west, align=left, font=\tiny, text=gray!55!black] at (-0.8,-3.5)
     {\textcolor{primary!55!black}{$\blacksquare$}~train \quad
      \textcolor{goldhl!60!black}{$\blacksquare$}~validate \quad
      \textcolor{gray!45}{$\blacksquare$}~not yet available};
\node[anchor=west, align=left, text width=13.0cm, font=\tiny, text=accent!75!black]
     at (-0.8,-4.1)
     {\textbf{Why the top row is fatal:} a model that has seen next month's traffic will
      predict this month's beautifully and be useless in production. The score is not
      merely optimistic --- it measures nothing that will ever happen again.};
\end{tikzpicture}}%
{Walk-forward (rolling-origin) validation. The training window only ever extends
backwards; each validation block sits strictly after its training data, exactly as
deployment will.}{walk-forward}

\begin{worked}[Choosing the Forecast Horizon]
A university wants to predict weekly VLE logins to size its servers.

\smallskip
\textbf{The question nobody asks first:} \emph{how far ahead, and why?} The horizon
is set by how long the resulting action takes, not by what the model can do.
\begin{tight}
  \item Auto-scaling cloud capacity: minutes of lead time needed $\Rightarrow$
        1-hour horizon.
  \item Provisioning extra servers: 2 weeks of procurement $\Rightarrow$ 2-week
        horizon.
  \item Budgeting next year's contract: 12-month horizon.
\end{tight}

\textbf{Why it changes the model.} At a 1-hour horizon, $y_{t-1}$ is available and
is overwhelmingly the best predictor --- a naive ``same as last hour'' baseline is
hard to beat. At a 12-month horizon, no recent lag is available at prediction time,
so the model must rely on trend, seasonality and enrolment forecasts, and error
will be several times larger.

\smallskip
\textbf{The consequence for evaluation.} A model quoted as ``8\% error'' is
meaningless without its horizon. Always report error \emph{per horizon}, and always
compare against the naive baseline at that same horizon --- which, at short
horizons, is a much stronger competitor than students expect.
\end{worked}

\section{Batch and Stream Processing}

\begin{definitionbox}[Batch versus Stream]
\term{Batch} processing runs over a bounded, stored dataset --- nightly, hourly.
\term{Stream} processing runs continuously over an unbounded sequence of events,
computing results incrementally as data arrives.
\end{definitionbox}

\dsfig{%
\begin{tikzpicture}[font=\scriptsize]
% tumbling
\node[font=\scriptsize\bfseries, text=primary, anchor=west] at (-0.6,2.35) {TUMBLING windows --- fixed, non-overlapping};
\foreach \i in {0,1,2,3}{
  \pgfmathsetmacro{\xx}{\i*2.6}
  \draw[rounded corners=2pt, draw=primary, fill=primary!10] (\xx,1.5) rectangle (\xx+2.4,2.1);
  \node[font=\tiny, text=primary!55!black] at (\xx+1.2,1.8) {0--5 min};}

% sliding
\node[font=\scriptsize\bfseries, text=greenhl!50!black, anchor=west] at (-0.6,0.95)
     {SLIDING windows --- fixed length, overlapping};
\foreach \i in {0,1,2,3,4,5}{
  \pgfmathsetmacro{\xx}{\i*1.3}
  \pgfmathsetmacro{\yy}{0.55-mod(\i,2)*0.35}
  \draw[rounded corners=2pt, draw=greenhl, fill=greenhl!10] (\xx,\yy) rectangle (\xx+2.4,\yy+0.3);}

% session
\node[font=\scriptsize\bfseries, text=goldhl!55!black, anchor=west] at (-0.6,-0.55)
     {SESSION windows --- bounded by a gap in activity};
\foreach \x/\w in {0/1.8, 2.9/3.2, 7.4/1.4, 9.8/2.6}{
  \draw[rounded corners=2pt, draw=goldhl, fill=goldhl!16] (\x,-1.35) rectangle (\x+\w,-0.95);}
\foreach \x in {2.1, 6.5, 9.1}{
  \node[font=\tiny, text=accent, anchor=center] at (\x,-1.15) {gap};}

\draw[-{Stealth[length=2mm]}, gray!55] (-0.6,-1.85) -- (12.8,-1.85);
\node[font=\tiny, text=gray!55!black] at (12.8,-2.15) {time};

\node[anchor=west, align=left, text width=13.0cm, font=\tiny, text=gray!55!black]
     at (-0.6,-2.75)
     {\textbf{Which to use:} tumbling for periodic reports (``failures per 5 minutes'');
      sliding for smooth, responsive alerting (``failed logins in the last 5 minutes,
      updated every 30 seconds''); session for user- or device-centric analysis where the
      natural unit is a burst of activity. Streaming systems also need a
      \textbf{watermark} --- a rule for how long to wait for late-arriving events before
      closing a window, which in IoT is unavoidable.};
\end{tikzpicture}}%
{Three windowing strategies for streaming data. Choosing the window is choosing what
``recent'' means, and it determines both detection latency and false-alarm
rate.}{windowing}

\section{Concept Drift}

\begin{definitionbox}[Concept Drift]
The relationship between features and target changes over time, so a model that was
correct becomes wrong without anything breaking. \term{Data drift} is a change in
the inputs' distribution; \term{concept drift} is a change in the mapping itself.
\end{definitionbox}

\dsfig{%
\begin{tikzpicture}
\begin{axis}[
  width=12.4cm, height=5.0cm,
  xlabel={\scriptsize weeks since deployment}, ylabel={\scriptsize model performance},
  xmin=0, xmax=52, ymin=0.35, ymax=0.95,
  axis lines=left, grid=major, grid style={gray!18},
  tick label style={font=\tiny}, label style={font=\tiny},
  legend style={font=\tiny, draw=gray!40, at={(0.99,0.99)}, anchor=north east},
]
% sudden drift
\addplot[accent, line width=1.2pt] coordinates {
 (0,0.86)(10,0.85)(19,0.845)(20,0.61)(30,0.60)(40,0.595)(52,0.59)};
\addlegendentry{sudden drift (a system change)}
% gradual drift
\addplot[goldhl, line width=1.2pt, domain=0:52, samples=60] {0.86 - 0.0055*x};
\addlegendentry{gradual drift (behaviour shifts)}
% with retraining
\addplot[greenhl, line width=1.2pt] coordinates {
 (0,0.86)(8,0.82)(8,0.86)(16,0.82)(16,0.86)(24,0.82)(24,0.86)(32,0.82)(32,0.86)
 (40,0.82)(40,0.86)(48,0.82)(48,0.86)(52,0.85)};
\addlegendentry{scheduled retraining}

\draw[gray!60, dashed] (axis cs:0,0.75) -- (axis cs:52,0.75);
\node[font=\tiny, text=gray!55!black, anchor=west] at (axis cs:0.5,0.72) {agreed minimum};
\draw[-{Stealth[length=1.8mm]}, accent] (axis cs:23,0.68) -- (axis cs:20.5,0.62);
\node[font=\tiny, accent, anchor=west] at (axis cs:23,0.69) {new VLE released};
\end{axis}
\end{tikzpicture}}%
{Three fates of a deployed model. Nothing about the code changed in any of them --- the
world did. Monitoring for this is the subject of \chref{mlops}.}{concept-drift}

\rowcolors{2}{white}{defbg}
\begin{longtable}{@{}p{3.0cm}p{4.8cm}p{6.4cm}@{}}
\toprule
\rowcolor{primary!15}
\textbf{Drift type} & \textbf{Looks like} & \textbf{Response} \\
\midrule
\endhead
Sudden & A cliff in the performance chart & Investigate the cause first --- usually an upstream change, not the model \\
Gradual & Slow decline over months & Scheduled retraining on a rolling window \\
Seasonal / recurring & Performance dips at the same point each cycle & Add the seasonal feature; do not retrain away a pattern you can model \\
Adversarial & Decline that accelerates after deployment & The attacker is adapting (\chref{cybersecurity}); retraining alone will not fix it \\
\bottomrule
\end{longtable}

%---------------------------------------------------------------------
\begin{fourlenses}{Time and Streams}
\lensLA{Engagement is strongly seasonal on a weekly and a semester cycle --- always
decompose before alerting, or you will raise an alarm every reading week. The
\emph{trajectory} of a student's residual is the signal: two consecutive weeks
below their own baseline is a far better trigger than any absolute threshold.
Validation must be walk-forward across cohorts, since a model trained on the 2025
cohort to predict 2024 proves nothing.}

\lensPM{Sprint data is a short, irregular series with strong autocorrelation ---
this sprint resembles the last. The naive baseline (``same as last sprint'') is
genuinely hard to beat, and saying so honestly is more valuable than a marginally
better model. Concept drift is guaranteed whenever the team composition changes,
which is why models here need a shorter memory than elsewhere.}

\lensIOT{This chapter is the technical bridge to \chref{iot}. Telemetry is the
purest streaming problem: unbounded, high-volume, out-of-order. Use sliding windows
for detection latency, tumbling for reporting, and set the watermark from the
worst-case network delay in your fleet. Sensors drift physically as well as
statistically --- a thermistor ages --- so drift monitoring must distinguish a
failing sensor from a changing process.}

\lensSEC{Windowed counts are the core feature type: failed logins per account per
5 minutes, distinct ports per host per minute. Diurnal and weekly seasonality is
strong, so a 3 a.m.\ spike matters and a 3 p.m.\ one may not --- decompose first.
Drift here is \emph{adversarial}, which makes it categorically different from the
other three lenses: the distribution is not shifting by accident, it is being
shifted deliberately, and scheduled retraining on recent data can be poisoned.}
\end{fourlenses}

%---------------------------------------------------------------------
\begin{lab}[Decompose, Feature, Validate]
\begin{lstlisting}[language=Python]
import pandas as pd, numpy as np
from statsmodels.tsa.seasonal import seasonal_decompose
from sklearn.model_selection import TimeSeriesSplit
from sklearn.ensemble import HistGradientBoostingRegressor
from sklearn.metrics import mean_absolute_error

s = df.set_index("date")["logins"].asfreq("W")

# --- 1. decompose before doing anything else --------------------------
seasonal_decompose(s, model="additive", period=12).plot()

# --- 2. features -- every one looks BACKWARDS only --------------------
X = pd.DataFrame(index=s.index)
for lag in (1, 2, 4, 12, 52):
    X[f"lag{lag}"] = s.shift(lag)                 # shift(+n) = past. Never shift(-n).
X["roll4_mean"] = s.shift(1).rolling(4).mean()    # shift(1) FIRST, then roll
X["roll4_std"]  = s.shift(1).rolling(4).std()
X["delta1"]     = s.shift(1) - s.shift(2)
X["week"]       = X.index.isocalendar().week.astype(int)
X["month"]      = X.index.month
y = s
X, y = X.dropna(), y.loc[X.dropna().index]

# --- 3. walk-forward validation, never random -------------------------
tscv = TimeSeriesSplit(n_splits=5)
naive_errs, model_errs = [], []
for tr, te in tscv.split(X):
    m = HistGradientBoostingRegressor().fit(X.iloc[tr], y.iloc[tr])
    model_errs.append(mean_absolute_error(y.iloc[te], m.predict(X.iloc[te])))
    naive_errs.append(mean_absolute_error(y.iloc[te], X.iloc[te]["lag1"]))

print(f"naive (last value) MAE : {np.mean(naive_errs):.2f}")
print(f"model MAE              : {np.mean(model_errs):.2f}")
# If the model does not clearly beat naive, the extra complexity is not earning its keep.

# --- 4. a simple drift monitor ----------------------------------------
recent, reference = model_errs[-1], np.mean(model_errs[:-1])
if recent > 1.5 * reference:
    print("DRIFT ALERT: recent error is 50% above the reference window")
\end{lstlisting}
\end{lab}

%---------------------------------------------------------------------
\begin{checkpoint}
\begin{tightnum}
  \item Why is \texttt{rolling(7, center=True)} a leakage bug in a forecasting
        feature?
  \item Traffic drops 40\% on 25 December. Is this an anomaly? What must you compute
        before answering?
  \item A model's accuracy fell from 0.87 to 0.62 in one week with no code change.
        What are the first two things you check?
\end{tightnum}
\end{checkpoint}

%---------------------------------------------------------------------
\begin{chaptersummary}
\begin{tight}
  \item Temporal observations are dependent, so the independence assumption behind
        Part III does not hold --- most importantly for validation.
  \item Decompose into trend, seasonality and residual; alert on the residual, not
        the raw value.
  \item Build lag, rolling, difference and calendar features, all strictly
        backward-looking.
  \item Validate walk-forward. Random k-fold on a time series trains on the future
        and produces meaningless scores.
  \item The forecast horizon is set by the lead time of the action it supports, and
        error must be reported per horizon against the naive baseline.
  \item Streams need windows --- tumbling, sliding or session --- plus a watermark
        for late data.
  \item Concept drift degrades deployed models silently. Monitor, and distinguish
        gradual, sudden, seasonal and adversarial drift.
\end{tight}
\end{chaptersummary}

\begin{reviewq}
  \item \textbf{(MCQ)} The correct validation scheme for a forecasting model is:
        (a)~random k-fold (b)~stratified k-fold (c)~walk-forward (d)~leave-one-out.
  \item \textbf{(MCQ)} Windows that are fixed-length and non-overlapping are called:
        (a)~sliding (b)~tumbling (c)~session (d)~expanding.
  \item \textbf{(Short)} Define concept drift and distinguish it from data drift.
  \item \textbf{(Short)} Why is the naive ``last value'' baseline so strong at short
        horizons, and what does that imply for how you report results?
  \item \textbf{(Short)} What is a watermark in stream processing and why does IoT
        need one?
  \item \textbf{(Applied)} You must alert on unusual failed-login volume per account.
        Specify the window type and length, the decomposition you would apply, and the
        threshold rule --- and justify each against the base-rate argument of
        \chref{probability}.
  \item \textbf{(Applied)} A deployed predictive-maintenance model degrades steadily
        over eight months. List four possible causes spanning sensors, process and
        model, and state the diagnostic that would distinguish them.
\end{reviewq}
